%%
%% Hybrid-Architecture Multimodal Large Models: A Survey
%% Built on the ACM acmart template (sample-manuscript.tex).
%%
\documentclass[manuscript,screen,review]{acmart}

\AtBeginDocument{%
  \providecommand\BibTeX{{%
    Bib\TeX}}}

\setcopyright{acmlicensed}
\copyrightyear{2026}
\acmYear{2026}
\acmDOI{XXXXXXX.XXXXXXX}
\acmConference[Conference acronym 'XX]{Make sure to enter the correct
  conference title from your rights confirmation email}{June 03--05,
  2026}{Woodstock, NY}
\acmISBN{978-1-4503-XXXX-X/2026/06}

\begin{document}

\title{Hybrid Visual Representations for Unified Multimodal Large Models:
A Survey of Architectures, Tokenizers, and Training Paradigms}

\author{Anonymous Author}
\affiliation{%
  \institution{Anonymous Institution}
  \city{Anonymous}
  \country{Anonymous}}
\email{anonymous@example.com}

\renewcommand{\shortauthors}{Anonymous}

\begin{abstract}
Unified multimodal large language models pursue a single backbone that
both understands and generates visual content, yet purely continuous and
purely discrete generation paradigms each impose well-known
limitations. Continuous representations preserve visual fidelity but
break the next-token prediction interface that scales language models;
discrete representations restore that interface but discard
high-frequency detail and weaken comprehension. Hybrid architectures
respond to this dilemma by mixing continuous and discrete signals along
the tokenizer, the input--output channels, or the prediction head. This
survey organizes the rapidly growing hybrid literature into five
strategies: embedding quantization, latent semantization, task-specific
decoupling, concatenated stacking, and fused merging. For each strategy,
we describe the central design idea, compare representative systems,
and discuss the trade-offs that shape current practice. We close with
open problems on representation conflict, training cost, and editing
fidelity that motivate future work.
\end{abstract}

\begin{CCSXML}
<ccs2012>
 <concept>
  <concept_id>10010147.10010178.10010224</concept_id>
  <concept_desc>Computing methodologies~Computer vision representations</concept_desc>
  <concept_significance>500</concept_significance>
 </concept>
 <concept>
  <concept_id>10010147.10010178.10010179.10010181</concept_id>
  <concept_desc>Computing methodologies~Natural language generation</concept_desc>
  <concept_significance>300</concept_significance>
 </concept>
</ccs2012>
\end{CCSXML}

\ccsdesc[500]{Computing methodologies~Computer vision representations}
\ccsdesc[300]{Computing methodologies~Natural language generation}

\keywords{Multimodal large language models, hybrid representation,
visual tokenizer, unified understanding and generation,
autoregression, diffusion}

\maketitle

\section{Introduction}

Unified multimodal large language models extend the next-token
prediction paradigm beyond text and ask a single backbone to read,
reason about, and synthesize images. Two paradigms dominate early work.
Continuous-token systems---such as Transfusion, OmniGen, and the
NExT-GPT family---feed the language model with dense visual embeddings
and decode pixels through diffusion or regression heads. Discrete-token
systems---such as Chameleon, Lumina-mGPT, and Emu3---quantize images
into a fixed vocabulary and reuse cross-entropy training without
modification. Both designs sit on opposite ends of an information
spectrum: continuous tokens carry rich detail at the cost of breaking
the symbolic interface, while discrete tokens preserve that interface
at the cost of representational bandwidth and semantic alignment.

Hybrid architectures occupy the middle of this spectrum. Rather than
choosing one signal type for the whole pipeline, they place continuous
and discrete components at different points of the model so that each
component handles the task it suits best. This survey defines a hybrid
multimodal large model as any system that, within one unified backbone,
combines continuous and discrete representations along at least one of
three axes: the visual tokenizer, the input--output interface, or the
prediction head. We organize this design space into five strategies.
\textit{Embedding quantization} converts continuous semantic embeddings
into discrete codes. \textit{Latent semantization} keeps the codes
continuous but injects high-level semantics into them.
\textit{Task-specific decoupling} routes understanding and generation
through separate visual paths. \textit{Concatenated stacking} feeds
both signal types to the backbone in parallel sequences.
\textit{Fused merging} combines the signals at the tokenizer level
through shared codebooks or hierarchical structures. The remainder of
this section explains why these strategies have emerged. Each
subsequent section develops one strategy in detail, ending with
comparative analysis and open problems.

\section{Embedding Quantization}

Embedding quantization compresses high-level continuous embeddings into
a discrete codebook, so that a single language model can predict visual
tokens with the standard next-token-prediction objective. The strategy
treats the image encoder, not the language model, as the locus of
hybridization: the encoder produces continuous semantic features, and
a vector-quantization stage discretizes these features before the
language model consumes them.

SEED pioneered this approach by reordering ViT patch features into a
one-dimensional causal sequence and training a codebook of 8192 entries
on top of a Causal Q-Former~\cite{ge2023seed}. The design choice is
deliberate: tokens with one-dimensional causal dependency match the
left-to-right attention of language models, and high-level semantics
prevent the codes from collapsing into pixel statistics. SEED-LLaMA
extends the tokenizer to a fully autoregressive multimodal model, while
SEED-X scales the recipe to interleaved image--text training.

Several follow-ups generalize the recipe across modalities. LaViT and
its video extension Video-LaViT discretize patches with a learned
selector that keeps only the informative tokens, improving efficiency
without losing semantic coverage. AnyGPT pushes the generality further
by reusing SEED for images, SpeechTokenizer for speech, and Encodec for
music, all funneled into one shared vocabulary~\cite{zhan2024anygpt}.
The unified vocabulary lets the language model treat new modalities as
new ``languages,'' so the architecture stays unchanged while the data
side absorbs the multimodal complexity.

The trade-off of embedding quantization is information loss at the
quantization step. Codebook collapse, limited vocabulary size, and the
single-stage joint optimization between reconstruction and semantic
alignment can hurt fine-grained perception. The next strategy responds
directly to this limitation.

\section{Latent Semantization}

Latent semantization reverses the direction of embedding quantization.
Instead of discretizing continuous semantic features, it keeps the
underlying tokens continuous and injects high-level semantics so that
one representation serves both understanding and generation. The
underlying claim is that the conflict between reconstruction and
alignment objectives is best resolved by disentangling them across
levels of abstraction, not by forcing a single quantization step to
satisfy both.

TokLIP demonstrates the principle by pairing an off-the-shelf
low-level VQ tokenizer (such as VQGAN) with a ViT-based token encoder
that distills CLIP-level semantics from the discrete codes. The
language model still trains on plain VQ tokens, while the token encoder
supplies the semantic signal needed for strong comprehension; the two
objectives never compete during tokenizer training~\cite{lin2025toklip}.
MMAR follows a similar principle by using a continuous masked
autoregressive head to fuse semantic and pixel information across
training stages. Show-o2 introduces a 3D causal VAE and a
spatial--temporal fusion module that concatenates SigLIP semantic
features with low-level structural features along the channel axis,
unifying image and video processing in one continuous space.
Ming-UniVision goes a step further with MingTok, a three-stage
continuous tokenizer in which a low-level encoder, a semantic decoder,
and a pixel decoder share one continuous latent space. The unified
space removes round-trips between heterogeneous representations and
yields more than a $3.5\times$ speedup in joint training convergence.
TUNA targets video by aligning a continuous tokenizer with both
temporal coherence and language grounding.

Latent semantization preserves visual fidelity because nothing is
quantized at the perceptual level, but it shifts the burden onto the
prediction head. The language model must now predict continuous
vectors, which calls for diffusion heads, flow matching, or
masked-regression objectives. These heads complicate training and
slow inference compared with pure cross-entropy. The next strategy
sidesteps this complication by using different paths for different
tasks.

\section{Task-Specific Decoupling}

Task-specific decoupling acknowledges that visual understanding and
visual generation make incompatible demands on a single representation.
Understanding favors abstract, text-aligned features that compress
appearance into semantics; generation favors pixel-faithful features
that retain texture and geometry. Rather than reconcile these demands
in one tokenizer, decoupled architectures route the two tasks through
separate visual encoders and a single shared backbone, so each task
receives the signal it needs while the backbone learns the joint
distribution.

Janus introduced the explicit decoupling principle: a SigLIP-style
encoder serves understanding, a VQ tokenizer serves generation, and a
unified transformer processes both streams~\cite{wu2024janus}.
JanusFlow integrates rectified flow into the same framework, so the
generation path inherits state-of-the-art continuous diffusion
modeling without abandoning the autoregressive backbone. Janus-Pro
scales the recipe to 7B parameters with optimized training and larger
data, narrowing the gap to specialist text-to-image models.

The decoupled design has been instantiated by many follow-ups, each
exploring a different combination of components.
\textbf{Autoregressive plus diffusion.} LMFusion, Mogao, BAGEL, and
OmniGen2 retain a transformer backbone that produces hidden states for
both text and conditioning, and delegate pixel synthesis to an external
diffusion decoder. BAGEL adopts a Mixture-of-Transformer-Experts design
in which understanding experts handle text and ViT tokens while
generation experts handle VAE tokens, sharing self-attention layers to
preserve cross-modal alignment. OmniGen2 keeps two separate decoding
paths with non-shared parameters and uses VAE features only inside the
diffusion decoder, which preserves the original multimodal
understanding ability of the language model.
\textbf{Visual autoregression.} VARGPT and VARGPTv1.1 replace the
diffusion decoder with a next-scale prediction head, generating images
through a coarse-to-fine sequence of discrete scales while keeping the
LLaVA-style understanding path intact. UniGen and UniGen1.5 adopt
similar scale-wise paradigms. Pisces and UniWorld-V1 explore lighter
combinations geared toward editing and in-context generation.
\textbf{Continuous-token autoregression.} UniFluid keeps the entire
sequence autoregressive but uses continuous visual tokens with a
diffusion-based prediction head, avoiding vector quantization
altogether and treating the loss balance between generation and
understanding as the central training problem.
\textbf{Lightweight task adapters.} Skywork UniPic, Skywork UniPic 2.0,
LightFusion, and Ovis-U1 contribute training-side innovations
including data-efficient stage curricula and parameter-efficient
adapters, allowing decoupled architectures to scale without
proportional compute growth.

Task-specific decoupling has become the dominant pattern of 2025
because it cleanly separates concerns. Its main cost is architectural
complexity: two encoders, a transformer, and often a diffusion or
scale-wise head must all be trained and tuned together, and the
information flow across these components can introduce mode-mixing
artifacts during interleaved generation.

\section{Concatenated Stacking}

Concatenated stacking takes a more conservative path. Instead of
designing custom encoders or prediction heads, it keeps off-the-shelf
semantic and pixel tokenizers, then concatenates their outputs into
one sequence that the language model consumes. The hybrid character
lives entirely in the input pipeline; the backbone remains a standard
transformer with a single language-modeling head.

ILLUME+ exemplifies this strategy. Its DualViTok runs two parallel
branches: a semantic branch built on a pre-trained text-aligned vision
encoder and a pixel branch that combines quantized features from the
semantic encoder with a CNN-based pixel encoder. The model adopts a
coarse-to-fine schedule, generating semantic tokens first and pixel
tokens second, and uses a continuous-input, discrete-output scheme to
prevent input-side information loss~\cite{huang2025illume}. A diffusion
decoder then refines the discrete outputs into high-resolution images.
The design preserves both texture and semantics while keeping the
language model interface uniform.

UniToken pursues a similar continuous-plus-discrete fusion at the
token-sequence level, demonstrating that the language model can absorb
heterogeneous visual streams when they are aligned in vocabulary
space. EMMA adds an emphasis on instruction tuning over the
concatenated input, showing that downstream prompt control benefits
from the explicit separation of semantic and pixel channels.

The strength of concatenated stacking is engineering simplicity:
existing tokenizers and language models slot together with minimal
architectural surgery. The cost is sequence length. Two parallel
streams either double the token count or multiply the vocabulary size,
which raises memory pressure during training and inference.

\section{Fused Merging}

Fused merging shares the dual-encoder intuition of concatenated
stacking but compresses the two streams back into one before they
reach the language model. The fusion may happen at the feature level,
through a shared codebook, or through a hierarchical codebook that
preserves the semantic--pixel relationship explicitly.

TokenFlow uses two encoders---a CLIP-style semantic encoder and a
VQGAN-style pixel encoder---and computes distances against two
codebooks in parallel. The combined distance selects a single code
index, which then retrieves quantized features from each codebook for
the respective decoders. The shared index keeps semantic and pixel
information consistent while still serving downstream understanding
and generation~\cite{qu2024tokenflow}. MUSE-VL adopts a closely related
recipe but concatenates the two streams along the dimension axis and
projects them through an MLP before quantization, which shortens the
sequence at the cost of a wider feature dimension.

SemHiTok takes the fusion idea one step further by introducing the
Semantic-Guided Hierarchical Codebook (SGHC). A pre-trained semantic
codebook organizes the top level, and a pixel sub-codebook attaches to
each semantic code. Quantization first selects a semantic code, then
selects the matching pixel code under it, and the two codes flatten
into one index for autoregressive prediction. The hierarchy decouples
the training of semantic and pixel objectives without losing the
ability to integrate into a standard next-token-prediction
backbone~\cite{chen2025semhitok}. The result is a unified tokenizer
that achieves competitive reconstruction fidelity and strong
multimodal understanding within one discrete interface.

Fused merging strikes a balance between the modular flexibility of
decoupled designs and the engineering simplicity of stacking
approaches. It avoids the doubled sequence length but requires careful
codebook design and stage-wise training to prevent the joint objective
from collapsing onto either side.

\section{Comparative Analysis}

The five hybrid strategies trade off along three axes: the location
of hybridization, the cost paid in training, and the quality
delivered at inference.
Embedding quantization places hybridization inside the tokenizer and
keeps the backbone untouched, which preserves the standard
cross-entropy training loop but caps fidelity at the quantization
ceiling. Latent semantization preserves fidelity by keeping tokens
continuous, but it forces the prediction head to leave the symbolic
interface. Task-specific decoupling separates the two tasks fully,
which yields the strongest empirical results in 2025 but multiplies
component count and training stages. Concatenated stacking minimizes
architectural change at the cost of sequence length, while fused
merging compresses two streams into one through codebook design at the
cost of additional engineering on the tokenizer side.

A second observation concerns the direction of recent work. Earlier
hybrid systems concentrated on the tokenizer (embedding quantization,
latent semantization), assuming that a better representation alone
would unify understanding and generation. More recent systems shift
the locus of hybridization to architecture-level decoupling and to
hierarchical codebooks, accepting that no single representation
satisfies both tasks and instead engineering the architecture to
acknowledge that conflict. The trend suggests that future hybrid
systems will continue to layer multiple strategies---for example,
combining a semanticized continuous tokenizer with task-specific heads
and a hierarchical codebook for editing---rather than choosing a
single design point.

A third observation concerns evaluation. The benchmarks used most often
in this literature, including MMBench, GenEval, and MJHQ, exercise
isolated capabilities. They do not yet test the property that hybrid
architectures advertise most prominently---multi-round, in-context
interleaving of understanding, generation, and editing within one
session. Until evaluation catches up, the relative merits of the five
strategies will remain partially obscured.

\section{Open Problems}

Three problems stand out for future work.

\textbf{Representation conflict.} Even the most carefully designed
hybrid representations still trade comprehension fidelity against
generation fidelity. Continuous tokenizers favor generation; discrete
tokenizers favor comprehension; hierarchical codebooks postpone the
conflict to a later stage. None of these choices eliminates the
underlying tension between abstraction and faithful reconstruction, and
a principled theory of when each strategy is preferable is still
missing.

\textbf{Training cost.} Hybrid systems typically demand more training
data and more stages than unimodal systems. The compute required to
train BAGEL, ILLUME+, or Ming-UniVision exceeds what most academic
groups can afford, which concentrates research on industrial labs and
risks narrowing the design space. Parameter-efficient finetuning,
distillation from large hybrids to smaller backbones, and curriculum
schedules tailored to hybrid signals remain underexplored directions.

\textbf{Editing fidelity.} The promise of unified models lies in
editing and in-context manipulation, not only in caption-then-generate
pipelines. Current hybrids handle text-to-image generation
competitively but lag specialist editors on instruction-driven
modification, color and identity preservation, and multi-round
refinement. Hierarchical codebooks and dual-encoder schemes have begun
to close the gap, yet substantial headroom remains, especially for
long-horizon, multi-turn editing where each round must respect the
previous round's pixels.

\section{Conclusion}

Hybrid multimodal large models occupy the design space between purely
continuous and purely discrete architectures. Five strategies have
emerged: embedding quantization, latent semantization, task-specific
decoupling, concatenated stacking, and fused merging. Each strategy
places the boundary between continuous and discrete signals at a
different point in the pipeline, and each pays a different price for
the unification it offers. The dominant pattern in 2025 favors
task-specific decoupling supplemented by hierarchical or fused
tokenizers, but the field has not converged. Open problems in
representation conflict, training cost, and editing fidelity will
shape the next generation of unified multimodal systems.

%% Bibliography
\bibliographystyle{ACM-Reference-Format}
\bibliography{hybrid-refs}

\end{document}
